

%\illustration{0.3}{image.jpg}{
%    Caption of the illustration (optional).
%    CC BY-SA 4.0 by X on \href{https://example.com/reference-to-image}{Y}
%}

% optionally define variables/limits for this problem



Your local jungle is being taken over by monkeys! Trees are quickly being
colonized by them! As a Brave Ape Pictures Collector, this is your chance of
taking sooo many pictures of primates that for sure will amaze your colleagues.

In particular, each day, one new monkey discovers your favourite tree containing
$n$ branches. From your
long experience observing primates, you know two things.
First, each branch of the tree only has room for a single monkey.
Second, monkeys are very social animals and always stick together in
groups, that is, the branches they occupy form a single connected component.

This particular invasive species of monkeys happens to be new to you and you
have not yet learned to distinguish them from each other. Still you wonder: how
many different pictures of the monkey colony could you take on each day, until
the tree is full of monkeys?

As an example, consider the first sample case, visualized in Figure~\ref{fig:monkeys}.
On the third day, there are four different pictures you can take of the monkey colony.

\begin{figure}[h]
    \centering
    \includegraphics[width=0.8\textwidth]{monkeys.drawio.pdf}
    \caption{
        Visualization of the first sample case on the third day.
        Branches are connected if they overlap
        (note the small gap between branches $1$ and $2$, and between branches $3$ and $4$).
    }
    \footnotesize
    Monkey image from \href{https://www.freevector.com/flat-brown-monkey-25484}{freevector.com}
    \label{fig:monkeys}
\end{figure}


Given the exact structure of your favourite tree, determine for each day from
$1$ to $n$ the number of different sets of branches the monkeys could occupy on
that day, modulo $10^9+7$.

